\documentclass[11pt,a4paper]{article}
\usepackage[margin=25mm,headheight=15pt]{geometry}
\usepackage[T1]{fontenc}
\usepackage[utf8]{inputenc}
\usepackage{lmodern,amsmath,amssymb,amsthm,mathtools,bm}
\usepackage{xcolor,graphicx,booktabs,tabularx,array,longtable}
\usepackage{microtype,enumitem,fancyhdr}
\usepackage{xurl}
\usepackage{hyperref}
\definecolor{accent}{HTML}{163B55}
\hypersetup{colorlinks=true,linkcolor=accent,citecolor=accent,urlcolor=accent,
 pdftitle={Generating the energy-constrained completion by completely positive maps},pdfauthor={Matthew J. Colbrook},
 pdfsubject={AIM 928: mathematical result and proof}}
\urlstyle{same}
\setlength{\parindent}{0pt}
\setlength{\parskip}{0.35em}
\setlength{\emergencystretch}{3em}
\setlist{itemsep=0.25em,topsep=0.3em}
\newtheorem{theorem}{Theorem}[section]
\newtheorem{lemma}[theorem]{Lemma}
\newtheorem{proposition}[theorem]{Proposition}
\newtheorem{corollary}[theorem]{Corollary}
\theoremstyle{remark}
\newtheorem{remark}[theorem]{Remark}
\newtheorem{definition}[theorem]{Definition}
\DeclareMathOperator{\Tr}{Tr}
\DeclareMathOperator{\supp}{supp}
\DeclareMathOperator{\Lip}{Lip}
\DeclareMathOperator{\Rea}{Re}
\newcommand{\R}{\mathbb R}
\newcommand{\C}{\mathbb C}
\newcommand{\E}{\mathbb E}
\newcommand{\N}{\mathbb N}
\newcommand{\ip}[2]{\langle #1,#2\rangle}
\pagestyle{fancy}
\fancyhf{}
\fancyhead[L]{\small\sffamily AIM 928}
\fancyhead[R]{\small\sffamily Matthew J. Colbrook}
\fancyfoot[C]{\thepage}
\renewcommand{\headrulewidth}{0.3pt}
\allowdisplaybreaks[2]
\title{\textbf{Generating the energy-constrained completion by completely positive maps}}
\author{Matthew J. Colbrook\\[3pt]
{\small Department of Applied Mathematics and Theoretical Physics}\\
{\small University of Cambridge, United Kingdom}\\
{\small\href{mailto:mjc249@maths.cam.ac.uk}{\texttt{mjc249@maths.cam.ac.uk}}}}
\date{3 November 2026}
\begin{document}
\maketitle
\begin{abstract}
We show that every element of the specified energy-constrained completion is a difference of two completely positive elements. A uniform decomposition in finite dimensions passes to the completion through a summable positive series. An exact infimum formula for the norm follows.
\end{abstract}
\noindent\textbf{Repository target.}
This manuscript concerns AIM problem 928 at revision
\texttt{5d7e26e8444e}. The precise scope is stated below.
\medskip

\section{Main result and proof}
\label{sec:928}

\paragraph{Answer.} The answer to the precise completion question is affirmative.
The key step is a dimension-independent completely positive decomposition on a
dense space of maps supported on finite input and output subspaces. An
absolutely convergent series then puts both components in the specified
closed cone, rather than merely in the larger cone of maps with finite
energy-constrained norm.

\begin{theorem}\label{thm:928}
Let $\mathcal H_A,\mathcal H_B$ be separable complex Hilbert spaces, and suppose
that $G\ge0$ on $\mathcal H_A$ has compact resolvent, minimum eigenvalue zero,
and eigenvalues tending to infinity. Fix $E>0$. Let $Y$ be the real vector
space of Hermiticity-preserving maps on trace class with finite ordinary
diamond norm, let $X_G$ be its completion for
$q(T)=\|T\|_{\diamond,E}^{G}$, and let $C_G$ be the closure in $X_G$ of its
completely positive cone. Then
\[
 X_G=C_G-C_G.
\]
For every $x\in X_G$ and $\eta>0$, the components may be chosen so that
$x=a-b$ and $q(a+b)\le 2q(x)+\eta$.
\end{theorem}

The quantitative constant two is sufficient for the theorem. A refinement
below replaces it by one, with the same arbitrarily small additive error.
No attainment of a minimal decomposition is asserted.

\subsection{Subnormalized inputs and finite-dimensional approximation}

Choose a unit vector $g\in\ker G$. If $\sigma\ge0$ is a subnormalized
input--ancilla state with $\operatorname{Tr}\sigma\le1$ and
$\operatorname{Tr}((G\otimes I)\sigma)\le E$, then
\[
 \|(T\otimes\mathrm{id})(\sigma)\|_1\le q(T).
 \tag{928.1}\label{eq:928-substate}
\]
Indeed, in an enlarged ancillary space add
$(1-\operatorname{Tr}\sigma)|g\rangle\langle g|$ in an orthogonal ancillary
flag. The resulting state is normalized and still has energy at most $E$.
Its output is a direct sum, whose trace norm is at least that of the output
on $\sigma$. An infinite-dimensional separable ancilla accommodates the
extra flag isometrically. This argument requires no positivity or
trace-preservation assumption on $T$.

Let $P_n$ be increasing finite-rank spectral projections of $G$, containing
the ground eigenspace, with $P_n\to I$ strongly. Choose $g_n\to\infty$ so that
$G\ge g_n(I-P_n)$. For every normalized state $\rho$ of input energy at most
$E$, the gentle projection bound gives
\[
 \|\rho-(P_n\otimes I)\rho(P_n\otimes I)\|_1
 \le 2\sqrt{\operatorname{Tr}((I-P_n)\rho_A)}
 \le 2\sqrt{E/g_n}.
\]
Consequently, for $T\in Y$,
\[
 q(T-T\circ\operatorname{Ad}_{P_n})
 \le 2\|T\|_\diamond\sqrt{E/g_n}\longrightarrow0,
 \qquad \operatorname{Ad}_V(\rho)=V\rho V^*.
 \tag{928.2}\label{eq:928-input}
\]

Fix $n$, put $d=\operatorname{rank}P_n$, and restrict the input to $P_n\mathcal
H_A\cong\mathbb C^d$. The resulting map has a trace-class Choi operator
\[
 J_n=\sum_{i,j=1}^d T(|e_i\rangle\langle e_j|)
                  \otimes|e_i\rangle\langle e_j|
 \in\mathcal T_1(\mathcal H_B\otimes\mathbb C^d).
\]
If $Q_m\to I$ strongly are finite-rank output projections, then
\[
 (Q_m\otimes I_d)J_n(Q_m\otimes I_d)\longrightarrow J_n
 \quad\hbox{in trace norm}.
\]
For any Hermiticity-preserving map $R$ with finite-dimensional input,
\[
 \|R\|_\diamond\le\|J_R\|_1.
 \tag{928.3}\label{eq:928-choi}
\]
To see this, purify an input state. Up to an ancillary isometry its output
has the form $(I\otimes\sqrt{\rho})J_R(I\otimes\sqrt{\rho})$, with
$\|\sqrt{\rho}\|\le1$; a transpose can be absorbed into the input convention.
Partial trace contracts the trace norm of Hermitian operators, so mixed
inputs cause no difficulty. Thus the maps
\[
 T_{n,m}(\rho)=Q_mT(P_n\rho P_n)Q_m
\]
approximate $T\circ\operatorname{Ad}_{P_n}$ in ordinary diamond norm, and
hence in $q$, as $m\to\infty$.

Let $D$ be the real linear space of maps with such finite input and output
support. Two of them can be placed in common larger subspaces, so $D$ is
indeed linear. Equations \eqref{eq:928-input}--\eqref{eq:928-choi} show that
$D$ is dense in $Y$ for $q$, and therefore dense in $X_G$.

\subsection{A uniform positive decomposition on the dense space}

Take $T\in D$, with input support $P=P_n$. On the finite-dimensional space
$P\mathcal H_A$, write $G_P=G|_{P\mathcal H_A}$ and set
\[
 W=(I+G_P/E)^{-1/2},\qquad S=T\circ\operatorname{Ad}_W.
\]
Here $W^{-1}$ is bounded because the input space is finite dimensional.
For any normalized state $\omega$ on this input space and an ancilla,
$\sigma=(W\otimes I)\omega(W\otimes I)$ has trace at most one and energy
at most $E$, since $WG_PW=E(I-W^2)\le EI$. By
\eqref{eq:928-substate},
\[
 \|S\|_\diamond\le q(T).
 \tag{928.4}\label{eq:928-weight}
\]

A finite-dimensional Hermiticity-preserving map admits completely positive
maps $S_+,S_-$ such that
\[
 S=S_+-S_-,\qquad
 \|S_++S_-\|_\diamond=\|S\|_\diamond.
 \tag{928.5}\label{eq:928-cpdecomp}
\]
For clarity, this is the Hermitian form of the diamond-norm semidefinite
program \cite{Watrous2009}. The precise Choi-form dual is given in
\cite[Theorem 6 and Section 3.2, pp.~10--12]{Watrous2012}.
For a Hermitian Choi matrix, average a feasible two-block dual with its
block swap. The two diagonal blocks then agree, and conjugation by the
two-by-two Hadamard matrix gives the constraints below.
If $J$ is the Hermitian Choi matrix, its dual
minimizes $t$ subject to $Z\ge J$, $Z\ge-J$, and
$\operatorname{Tr}_B Z\le tI$. Strong duality and attainment hold in finite
dimensions. The Choi matrices $(Z+J)/2$ and $(Z-J)/2$ give
\eqref{eq:928-cpdecomp}; a completely positive map has diamond norm equal to
the operator norm of its trace functional.

Extend $S_\pm$ by zero outside $P\mathcal H_A$ and define ordinary bounded
completely positive maps
\[
 T_\pm=S_\pm\circ\operatorname{Ad}_{W^{-1}P}.
\]
They have finite input and output support, and $T=T_+-T_-$. For any full
input--ancilla state $\rho$ of energy at most $E$, put
\[
 \sigma=(P\otimes I)\rho(P\otimes I),\qquad
 \tau=(W^{-1}\otimes I)\sigma(W^{-1}\otimes I).
\]
All operators in this formula act on the finite input subspace, and
\[
 \operatorname{Tr}\tau
 =\operatorname{Tr}\sigma+E^{-1}\operatorname{Tr}((G_P\otimes I)\sigma)
 \le2.
\]
Moreover, $(T_++T_-)\otimes\mathrm{id}(\rho)
=(S_++S_-)\otimes\mathrm{id}(\tau)$. Therefore
\[
 q(T_++T_-)\le2\|S_++S_-\|_\diamond\le2q(T).
 \tag{928.6}\label{eq:928-uniform}
\]
Complete positivity also gives $q(T_\pm)\le q(T_++T_-)$, because the
positive output of either component has trace at most the trace of their
sum. This establishes the required uniform decomposition without using an
infinite-dimensional normal-map decomposition or an unbounded inverse.

\subsection{An absolutely convergent series proves exact cone membership}

Let $x\in X_G$ and $\epsilon>0$. Choose $x_k\in D$ with
$q(x-x_k)\le\epsilon2^{-k}$ for $k\ge1$. Put $d_1=x_1$ and
$d_k=x_k-x_{k-1}$ for $k\ge2$. Then
\[
 \sum_{k\ge1}q(d_k)
 \le q(x)+2\sum_{k\ge1}q(x-x_k)
 \le q(x)+2\epsilon.
 \tag{928.7}\label{eq:928-series}
\]
For each $k$, write $d_k=a_k-b_k$ using
\eqref{eq:928-uniform}. The series $a=\sum_ka_k$ and $b=\sum_kb_k$
converge in $X_G$, since $q(a_k),q(b_k)\le2q(d_k)$.
Every finite partial sum is an ordinary bounded completely positive map.
By the definition of $C_G$, both limits belong to $C_G$. Finally,
\[
 a-b=\sum_{k\ge1}d_k=\lim_{k\to\infty}x_k=x,
 \qquad q(a+b)\le2q(x)+4\epsilon.
\]
Taking $\epsilon<\eta/4$ proves Theorem~\ref{thm:928}.

\subsection{Refinement of the quantitative bound}

In fact, for $T\in D$ there is a decomposition with
$q(T_++T_-)=q(T)$. Here is a direct finite-dimensional derivation.
Use an eigenbasis of $G_P$ and let $J=J_T$. Purification and trace-norm
duality identify $q(T)$ with
\[
 \begin{aligned}
 \text{maximize}\quad&\operatorname{Tr}(JX),\\
 \text{subject to}\quad&\rho\ge0,\quad\operatorname{Tr}\rho=1,
   \quad\operatorname{Tr}(G_P\rho)\le E,\\
 &-I_B\otimes\rho\le X\le I_B\otimes\rho.
 \end{aligned}
 \tag{928.8}\label{eq:928-primal}
\]
The input transpose in the Choi convention leaves $G_P$ unchanged.
The finite-input norm equals the full-space norm of $T$: any compressed
full-space input is subnormalized and can be padded within the input
subspace, as in \eqref{eq:928-substate}.

Let $A,B\ge0$ multiply the two matrix inequalities in
\eqref{eq:928-primal}, and let $\lambda\ge0$ and $\mu\in\mathbb R$
multiply the energy and trace constraints. Supremizing the Lagrangian
forces $A-B=J$ and
$\operatorname{Tr}_B(A+B)\le\mu I+\lambda G_P$. Thus, with $Z=A+B$, the dual is
\[
 \begin{aligned}
 \text{minimize}\quad&\mu+\lambda E,\\
 \text{subject to}\quad&Z\ge J,\quad Z\ge-J,\\
 &\operatorname{Tr}_B Z\le\mu I+\lambda G_P,\qquad\lambda\ge0.
 \end{aligned}
 \tag{928.9}\label{eq:928-dual}
\]
There is a positive-definite density matrix of energy strictly below $E$,
obtained by concentrating sufficiently strongly on a ground vector. With
$X=0$ this gives strict primal feasibility. Hence the dual optimum is
attained and equals $q(T)$.

The positive Choi matrices $(Z\pm J)/2$ give $T_\pm$. Their sum has
energy-constrained norm at most $\mu+\lambda E=q(T)$; the reverse inequality
follows by applying the triangle inequality to the two positive outputs.
Extending outside the input subspace is harmless: the dual inequality on a
ground vector implies $\mu\ge0$, so it remains valid on the omitted spectral
subspace. Substituting this improved decomposition into
\eqref{eq:928-series} proves that, for every $\eta>0$, one can arrange
\[
 x=a-b,\quad a,b\in C_G,\qquad q(a+b)\le q(x)+\eta.
\]

\begin{corollary}[Exact infimum formula]\label{cor:928-infnorm}
For every $x\in X_G$,
\[
 q(x)=\inf\{q(a+b):a,b\in C_G,\ x=a-b\}.
\]
The formula concerns an infimum; it does not assert that an optimal pair exists.
\end{corollary}
\begin{proof}
For ordinary bounded completely positive maps $A,B$, positivity of the two
outputs on every admissible input gives
\[
 \|(A-B)\otimes\mathrm{id}(\rho)\|_1
 \le\operatorname{Tr}((A+B)\otimes\mathrm{id}(\rho)).
\]
Taking the supremum yields $q(A-B)\le q(A+B)$. Given $a,b\in C_G$, choose
ordinary bounded completely positive maps $A_n\to a$ and $B_n\to b$ in
$q$. Continuity of the norm passes this inequality to the limit, giving
$q(a-b)\le q(a+b)$. This proves that the displayed infimum is at least
$q(x)$. The decompositions with $q(a+b)\le q(x)+\eta$ established above,
for arbitrary $\eta>0$, prove the converse inequality.
\end{proof}

\paragraph{Relation to the cited question.}
Shirokov's conjecture in Section~5 of \cite{Shirokov2019} asks for components
in the completed cone $\mathfrak F_G^0$. The decomposition following
Lemma~2.19 in \cite{vanLuijk2025} produces components in the larger finite
energy-norm cone. The series above addresses precisely the extra closure
requirement: its two components are norm limits of ordinary bounded
completely positive maps. This establishes the catalogue formulation
\cite{AIM} without identifying those two cones.


\begin{thebibliography}{99}
\small\raggedright
\setlength{\itemsep}{3pt}

\bibitem{Watrous2009}
J.~Watrous,
\emph{Semidefinite programs for completely bounded norms},
Theory Comput. \textbf{5} (2009), 217--238.
\href{https://doi.org/10.4086/toc.2009.v005a011}{doi:10.4086/toc.2009.v005a011}.

\bibitem{Shirokov2019}
M.~E.~Shirokov,
\emph{On completion of the cone of completely positive linear maps with
respect to the energy-constrained diamond norm},
Lobachevskii J. Math. \textbf{40} (2019), 1549--1568.
\href{https://doi.org/10.1134/S1995080219100238}{doi:10.1134/S1995080219100238}.
\href{https://arxiv.org/abs/1810.10922}{arXiv:1810.10922}, Section 5.

\bibitem{Watrous2012}
J.~Watrous,
\emph{Simpler semidefinite programs for completely bounded norms},
\href{https://arxiv.org/abs/1207.5726v2}{arXiv:1207.5726v2},
2 August 2012, Theorem 6, p.~10, and Section 3.2, pp.~11--12.

\bibitem{vanLuijk2025}
L.~van Luijk,
\emph{Energy-Limited Quantum Dynamics},
Comm. Math. Phys. \textbf{406} (2025), article 120.
\href{https://doi.org/10.1007/s00220-025-05282-w}{doi:10.1007/s00220-025-05282-w}.
Lemma 2.19 and footnote 11.

\bibitem{AIM}
MColbrook/AIM, \emph{AIM 928: Does the closed completely positive cone generate the energy-constrained completion?},
repository statement at revision \texttt{5d7e26e8444e}.
\url{https://github.com/MColbrook/AIM/blob/5d7e26e8444eb98de91988454564e45e011107c4/problems/928-does-the-closed-completely-positive-cone-generate-the-energy-constrained-completion.md}.
\end{thebibliography}
\end{document}
